\documentclass[../../main.tex]{subfiles}

\begin{document}

\section{开集、闭集以及收敛序列}\label{section:ch09-开集闭集收敛序列}

在欧氏空间和一般赋范线性空间中研究的许多概念可自然地、有用地推广到一般的度量空间. 它们不依赖于线性结构.

\begin{definition}[开球与开集]\label{definition:ch09-开球开集}
令 $(X,\rho)$ 为度量空间. 对 $X$ 中的点 $x$ 和 $r>0$, 集合
\begin{align*}
B(x,r)=\{x'\in X\mid\rho(x',x)<r\}
\end{align*}
称为中心在 $x$、半径为 $r$ 的开球. $X$ 的子集 $\mathcal{O}$ 称为开的, 若对每个点 $x\in\mathcal{O}$, 存在中心在 $x$ 包含于 $\mathcal{O}$ 的开球. 对于点 $x\in X$, 包含 $x$ 的开集称为 $x$ 的邻域.
\end{definition}

我们必须验证这里的定义是一致的, 即开球是开的. 根据开集的定义, 要证明 $B(x,r)$ 是开的, 仅须证明: 若 $x'\in B(x,r)$ 且 $r'=r-\rho(x',x)$, 则 $B(x',r')\subseteq B(x,r)$. 为验证这一点, 令 $y\in B(x',r')$, 则 $\rho(x',y)<r'$, 因此, 根据三角不等式,
\begin{align*}
\rho(y,x)\leqslant\rho(y,x')+\rho(x',x)<r'+\rho(x',x)=r.
\end{align*}
因此 $B(x',r')\subseteq B(x,r)$.

\begin{proposition}\label{proposition:ch09-开集运算}
令 $X$ 为度量空间. 全空间 $X$ 和空集 $\emptyset$ 是开的, $X$ 的任何两个开子集的交是开的, $X$ 的任何开子集族的并是开的.
\end{proposition}

\begin{proof}
显然 $X$ 和 $\emptyset$ 是开的且开集族的并是开的. 令 $\mathcal{O}_{1}$ 和 $\mathcal{O}_{2}$ 为 $X$ 的开子集. 若这两个集合不交, 则它们的交集是空集, 它自然是开的. 否则, 令 $x$ 属于 $\mathcal{O}_{1}\bigcap\mathcal{O}_{2}$. 由于 $\mathcal{O}_{1}$ 和 $\mathcal{O}_{2}$ 是包含 $x$ 的开集, 存在正数 $\delta_{1}$ 和 $\delta_{2}$ 使得 $B(x,\delta_{1})\subseteq\mathcal{O}_{1}$ 且 $B(x,\delta_{2})\subseteq\mathcal{O}_{2}$. 定义 $\delta=\min\{\delta_{1},\delta_{2}\}$. 则开球 $B(x,\delta)$ 包含于 $\mathcal{O}_{1}\bigcap\mathcal{O}_{2}$. 因此 $\mathcal{O}_{1}\bigcap\mathcal{O}_{2}$ 是开的.
\end{proof}

以下命题的证明留作练习, 在度量空间 $X$ 是度量空间 $Y$ 的子空间的情形下, 提供了用 $Y$ 的开子集刻画 $X$ 的开子集的描述.

\begin{proposition}\label{proposition:ch09-子空间开集}
令 $X$ 为度量空间 $Y$ 的子空间, 而 $E$ 是 $X$ 的子集. 则 $E$ 在 $X$ 中是开的当且仅当 $E=X\bigcap\mathcal{O}$, 其中 $\mathcal{O}$ 在 $Y$ 中是开的.
\end{proposition}

\begin{definition}[闭包点与闭包]\label{definition:ch09-闭包点闭包}
对于度量空间 $X$ 的子集 $E$, 点 $x\in X$ 称为 $E$ 的闭包点, 若 $x$ 的每个邻域至少包含 $E$ 中的一个点. $E$ 的全体闭包点称为 $E$ 的闭包且记为 $\overline{E}$.
\end{definition}

显然我们总是有 $E\subseteq\overline{E}$. 若 $E$ 包含它的所有闭包点, 即 $E=\overline{E}$, 则称集合 $E$ 为闭的. 对于度量空间 $(X,\rho)$ 的点 $x$ 和 $r>0$, 集合 $\overline{B}(x,r)\equiv\{x'\in X\mid\rho(x',x)\leqslant r\}$ 称为中心在 $x$ 半径为 $r$ 的闭球. 从度量的三角不等式得出 $\overline{B}(x,r)$ 是包含 $B(x,r)$ 的闭集. 在赋范线性空间 $X$ 中, 我们称 $B(0,1)$ 为开单位球, 而 $\overline{B}(0,1)$ 为闭单位球.

\begin{proposition}\label{proposition:ch09-闭包最小闭集}
对于度量空间 $X$ 的子集 $E$, 它的闭包 $\overline{E}$ 是闭的. 此外, $\overline{E}$ 在以下意义下是 $X$ 的包含 $E$ 的最小闭子集: 若 $F$ 是闭集且 $E\subseteq F$, 则 $\overline{E}\subseteq F$.
\end{proposition}

\begin{proof}
集合 $\overline{E}$ 是闭的, 若它包含它的所有闭包点. 令 $x$ 为 $\overline{E}$ 的闭包点. 考虑 $x$ 的邻域 $\mathcal{U}_{x}$. 存在点 $x'\in\overline{E}\bigcap\mathcal{U}_{x}$. 由于 $x'$ 是 $E$ 的闭包点, 而 $\mathcal{U}_{x}$ 是 $x'$ 的一个邻域, 存在点 $x''\in E\bigcap\mathcal{U}_{x}$. 因此 $x$ 的每个邻域包含 $E$ 的点, 从而 $x\in\overline{E}$. 因此集合 $\overline{E}$ 是闭的. 显然, 若 $A\subseteq B$, 则 $\overline{A}\subseteq\overline{B}$, 因此若 $F$ 是闭的且包含 $E$, 则 $\overline{E}\subseteq F=F$.
\end{proof}

\begin{proposition}\label{proposition:ch09-开集补集闭集}
度量空间 $X$ 的子集是开的当且仅当它在 $X$ 中的补是闭的.
\end{proposition}

\begin{proof}
首先假定 $E$ 在 $X$ 中是开的. 令 $x$ 为 $X\sim E$ 的闭包的点. 则 $x$ 不属于 $E$, 这是因为, 否则将有 $x$ 的邻域包含于 $E$, 因此与 $X\sim E$ 不交. 因此 $x$ 属于 $X\sim E$, 从而 $X\sim E$ 是闭的. 现在假定 $X\sim E$ 是闭的. 令 $x$ 属于 $E$. 则必有 $x$ 的邻域包含于 $E$, 否则 $x$ 的每个邻域包含 $X\sim E$ 的点, 因此 $x$ 是 $X\sim E$ 的闭包的点. 由于 $X\sim E$ 是闭的, $x$ 也属于 $X\sim E$. 这是一个矛盾.
\end{proof}

由于 $X\sim[X\sim E]=E$, 从前一个命题得出一个集合是闭的当且仅当它的补是开的. 因此, 根据 De Morgan 等式, 命题 1 可通过闭集重述如下.

\begin{proposition}\label{proposition:ch09-闭集运算}
令 $X$ 为度量空间. 全空间 $X$ 和空集 $\emptyset$ 是闭的, $X$ 的任何两个闭子集的并是闭的, $X$ 的任何闭子集族的交是闭的.
\end{proposition}

我们已经定义了赋范线性空间的序列收敛于一点意味着什么. 以下是收敛性在度量空间的自然推广.

\begin{definition}[序列收敛]\label{definition:ch09-序列收敛}
度量空间 $(X,\rho)$ 的序列 $\{x_{n}\}$ 称为收敛到点 $x\in X$, 若
\begin{align*}
\lim_{n\to\infty}\rho(x_{n},x)=0.
\end{align*}
即对每个 $\varepsilon>0$, 存在指标 $N$ 使得对每个 $n\geqslant N$, $\rho(x_{n},x)<\varepsilon$. 序列收敛的点 $x$ 称为序列的极限. 我们常以 $\{x_{n}\}\to x$ 表示 $\{x_{n}\}$ 收敛于 $x$.
\end{definition}

度量空间中的序列至多收敛到一个点. 事实上, 给定度量空间 $X$ 中的两点 $u$, $v$, 设 $r=\rho(u,v)/2$. 我们从度量 $\rho$ 的三角不等式推出球 $B(u,r)$ 和 $B(v,r)$ 是不交的. 因此一个序列不可能既收敛到 $u$ 又收敛到 $v$. 此外, 收敛性可重述如下: $\{x_{n}\}$ 收敛到极限 $x$, 若对 $x$ 的任何邻域 $\mathcal{O}$, 除有限项外序列的所有项都属于 $\mathcal{O}$. 自然地, 对于 $X$ 的子集 $E$ 和使得对所有 $n$, $x_{n}$ 属于 $E$ 的序列 $\{x_{n}\}$, 我们说 $\{x_{n}\}$ 是 $E$ 中的序列.

\begin{proposition}\label{proposition:ch09-闭包点序列极限}
对于度量空间 $X$ 的子集 $E$, 点 $x\in E$ 是 $E$ 的闭包点当且仅当 $x$ 是 $E$ 中的序列的极限. 因此 $E$ 是闭的, 当且仅当每当 $E$ 的序列收敛到极限 $x\in X$, 该极限 $x$ 属于 $E$.
\end{proposition}

\begin{proof}
仅需证明第一个断言. 首先假定 $x$ 属于 $\overline{E}$. 对每个自然数 $n$, 由于 $B(x,1/n)\bigcap E\neq\emptyset$, 我们可以选一个点, 将它记为 $x_{n}$, 它属于 $B(x,1/n)\bigcap E$. 则 $\{x_{n}\}$ 是 $E$ 中的序列, 且我们断言它收敛到 $x$. 的确, 令 $\varepsilon>0$. 选取指标 $N$ 满足 $1/N<\varepsilon$. 则若 $n\geqslant N$,
\begin{align*}
\rho(x_{n},x)<1/n<1/N<\varepsilon.
\end{align*}
因此 $\{x_{n}\}$ 收敛到 $x$. 反过来, 若 $E$ 中的序列收敛到 $x$, 则每个中心在 $x$ 的球包含该序列的无限多个点, 因此包含 $E$ 中的点. 因此 $x\in\overline{E}$.
\end{proof}

一般来说, 集合上的度量的改变将会改变一个集合是开的含义, 从而改变一个集合是闭的含义. 序列的收敛的含义也会改变. 例如, 对于集合 $X$ 上的离散度量, 每个子集是开的, 每个子集是闭的, 且一个序列收敛到一个极限当且仅当除有限项外该序列的所有项都等于该极限. 以下命题的证明留作练习, 它告诉我们对于集合上的等价度量, 开集是相同的, 因此闭集是相同的且序列的收敛是相同的.

\begin{proposition}\label{proposition:ch09-等价度量开集相同}
令 $\rho$ 和 $\sigma$ 为非空集合 $X$ 上的等价度量. 则 $X$ 的子集在度量空间 $(X,\rho)$ 中是开的当且仅当它在度量空间 $(X,\sigma)$ 中是开的.
\end{proposition}

\begin{exercise}\label{exercise:ch09-13}
在一个度量空间 $X$ 中, 对于 $r>0$ 以及 $X$ 中的两个不同点 $u$ 和 $v$, 是否可能有 $B(u,r)=B(v,r)$? 在欧氏空间 $\mathbf{R}^{n}$ 中是否可能? 在赋范线性空间中是否可能?
\end{exercise}

\begin{exercise}\label{exercise:ch09-14}
令 $(X,\rho)$ 为度量空间使得 $\{u_{n}\}\to u$ 且 $\{v_{n}\}\to v$. 证明 $\{\rho(u_{n},v_{n})\}\to\rho(u,v)$.
\end{exercise}

\begin{exercise}\label{exercise:ch09-15}
令 $X$ 为度量空间, $x$ 属于 $X$ 而 $r>0$.
\begin{enumerate}[(i)]
\item 证明 $\overline{B}(x,r)$ 是闭的且包含 $B(x,r)$.
\item 证明在赋范线性空间 $X$ 中闭球 $\overline{B}(x,r)$ 是开球 $B(x,r)$ 的闭包, 但在一般的度量空间中这不成立.
\end{enumerate}
\end{exercise}

\begin{exercise}\label{exercise:ch09-16}
证明命题 2.
\end{exercise}

\begin{exercise}\label{exercise:ch09-17}
证明命题 7.
\end{exercise}

\begin{exercise}\label{exercise:ch09-18}
令 $X$ 为度量空间 $Y$ 的子空间而 $A$ 是 $X$ 的子集. 证明 $A$ 在 $X$ 中是闭的当且仅当 $A=X\bigcap F$, 其中 $F$ 在 $Y$ 中是闭的.
\end{exercise}

\begin{exercise}\label{exercise:ch09-19}
令 $X$ 为度量空间 $Y$ 的子空间.
\begin{enumerate}[(i)]
\item 若 $\mathcal{O}$ 是度量空间 $X$ 的开子集, 它是 $Y$ 的开子集吗? 若 $X$ 是 $Y$ 的开子集呢?
\item 若 $F$ 是度量空间 $X$ 的闭子集, 它是 $Y$ 的闭子集吗? 若 $X$ 是 $Y$ 的闭子集呢?
\end{enumerate}
\end{exercise}

\begin{exercise}\label{exercise:ch09-20}
对于度量空间 $X$ 的子集 $E$, 点 $x\in X$ 称为 $E$ 的内点, 若存在中心在 $x$ 包含于 $E$ 的开球; $E$ 的全体内点称为 $E$ 的内部且记为 $\mathrm{int}E$. 证明 $\mathrm{int}E$ 总是开的, 而 $E$ 是开的当且仅当 $\mathrm{int}E=E$.
\end{exercise}

\begin{exercise}\label{exercise:ch09-21}
对度量空间 $X$ 的子集 $E$, 点 $x\in X$ 称为 $E$ 的外部点, 若存在中心在 $x$ 包含于 $X\sim E$ 的开球; $E$ 的全体外部点称为 $E$ 的外部且记为 $\mathrm{ext}E$. 证明 $\mathrm{ext}E$ 总是开的, $E$ 是闭的当且仅当 $X\sim E=\mathrm{ext}E$.
\end{exercise}

\begin{exercise}\label{exercise:ch09-22}
对于度量空间 $X$ 的子集 $E$, 点 $x\in X$ 称为 $E$ 的边界点, 若每个中心在 $x$ 的开球包含 $E$ 与 $X\sim E$ 的点; $E$ 的全体边界点称为 $E$ 的边界且记为 $\mathrm{bd}E$. 证明: (i) $\mathrm{bd}E$ 总是闭的, (ii) $E$ 是开的当且仅当 $E\bigcap\mathrm{bd}E=\emptyset$, (iii) $E$ 是闭的当且仅当 $\mathrm{bd}E\subseteq E$.
\end{exercise}

\begin{exercise}\label{exercise:ch09-23}
令 $A$ 和 $B$ 为度量空间 $X$ 的子集. 证明: 若 $A\subseteq B$, 则 $\overline{A}\subseteq\overline{B}$. 证明 $\overline{A\bigcup B)}=\overline{A}\bigcup\overline{B}$ 且 $\overline{A\bigcap B)}\subseteq\overline{A}\bigcap\overline{B}$.
\end{exercise}

\begin{exercise}\label{exercise:ch09-24}
证明对度量空间 $X$ 的子集 $E$, $E$ 的闭包是 $X$ 的所有包含 $E$ 的闭子集的交.
\end{exercise}

\begin{exercise}\label{exercise:ch09-25}
令 $\rho$ 为集合 $X$ 上的度量. 对所有 $u$, $v\in X$, 定义
\begin{align*}
\tau(u,v)=\frac{\rho(u,v)}{1+\rho(u,v)}.
\end{align*}
验证 $\tau$ 是 $X$ 上的有界度量, 且序列关于度量 $\rho$ 的收敛性与关于度量 $\tau$ 的收敛性是相同的. 得出以下结论: 关于度量 $\rho$ 是闭的集合关于度量 $\tau$ 是闭的, 且关于度量 $\rho$ 是开的集合关于度量 $\tau$ 是开的. 度量 $\rho$ 和 $\tau$ 是等价的吗?
\end{exercise}

\end{document}